% This chapter presents the current state of the art in the domain and talks about other similar work that has been done in this area. It also establishes the novelty of our work by highlighting the differences between the existing work and our work.

% We will keep updating this chapter (especially if our project is research-intensive) as our research proceeds and we come across more work related to our problem.

% Of course, we take inspiration from \cite{einstein} but wish the work was typeset in \LaTeX \cite{knuthwebsite}, e.g. by taking help from \cite{latexcompanion}.

\section{Introduction}

The rapid advancement of artificial intelligence (AI), particularly large language models (LLMs), has significantly reshaped organizational knowledge management and service delivery. Retrieval-Augmented Generation (RAG) enhances LLMs by grounding responses in external knowledge sources, improving factual accuracy, reducing hallucinations, and increasing transparency \cite{lewis2020rag,gao2023ragsurvey,fan2024ragllm}. In parallel, nonprofit and humanitarian organizations are increasingly adopting conversational systems for service navigation, donor engagement, and community outreach \cite{folstad2018socialgood,siddiqi2024bablibot,munir2021bablibot}.

This literature review synthesizes research across four interconnected areas to inform the design of a multi-domain organizational assistant for a Pakistani NGO. First, it examines RAG-based conversational systems, including advances in persona adaptation, multi-turn dialogue, and multimodal interaction. Second, it reviews chatbot deployments for social good, with a focus on low- and middle-income country (LMIC) contexts and local-language systems in Urdu and Roman Urdu. Third, it explores system-level considerations, including human-in-the-loop escalation, operational dashboards, and analytics. Finally, it considers infrastructure, security, and deployment practices required for scalable and reliable systems.

A key gap emerges from this synthesis: existing chatbot systems in Pakistan are predominantly single-domain and loosely integrated with organizational workflows. There is limited evidence of a unified, multi-domain RAG-based assistant capable of supporting donations, healthcare services, blood programs, and welfare initiatives across English, Urdu, and Roman Urdu. The proposed Alkhidmat Public Chat Portal addresses this gap by combining RAG grounding, multilingual interaction, multimodal inputs, confidence-aware escalation, and integrated operational tooling within a single system.

\section{State of the Art in RAG-Based Organizational Assistants}
For tasks that require a lot of knowledge, Retrieval-Augmented Generation (RAG) was introduced as a general framework that combines pretrained sequence to sequence models with a non-parametric memory of retrieved documents \cite{lewis2020rag}. Lewis et al. note that models augmented with retrieval not only generate more accurate and factual responses but also outperform both purely parametric models and traditional retrieve-and-extract approaches in open-domain question answering.

The RAG framework has since been organised and made clearer by a number of survey studies.  A thorough analysis of RAG in large language models is given by Gao et al. \cite{gao2023ragsurvey}, who also discuss naive, advanced, and modular architectures and explain how the retrieval, generation, and augmentation components interact.  Similar to this, Fan et al. concentrate on Retrieval-Augmented Large Language Models (RA-LLMs), looking at their architectures, training methods, and uses in domains like recommendation systems and customer service \cite{fan2024ragllm}.  When taken as a whole, these surveys show that RAG has emerged as a crucial strategy for producing organisational assistants who are trustworthy and able to maintain a solid foundation in current, authoritative knowledge.

At the same time, research on conversational AI has made significant progress in areas such as persona modeling, multi-turn context tracking, and socially aware interactions. Surveys of neural conversational agents show a clear shift from simple, single-turn question answering toward multi-turn, task-oriented, and open-domain dialogue systems that can adapt both style and content to meet users’ needs \cite{kim2023convai}. Persona-aware dialogue models take this a step further, adjusting tone and content based on user or role profiles—an important consideration when designing assistants for beneficiaries, donors, or staff who may have different information requirements \cite{singh2024persona}.

From these trends, contemporary organizational assistants typically:
\begin{itemize}\setlength{\itemsep}{0pt}
    \item use RAG to ground answers in internal knowledge bases;
    \item support multi-channel access (web, mobile, messaging apps);
    \item adapt to different user personas and roles;
    \item include monitoring, guardrails, and human-in-the-loop escalation.
\end{itemize}
The proposed Alkhidmat assistant follows this pattern by coupling an LLM front-end with domain-structured knowledge on donations, healthcare, blood services, and welfare programmes, plus escalation and analytics for staff.

\section{AI Chatbots in Nonprofits and Civil Society}
While RAG provides the technical foundation for knowledge-grounded systems, its application in real-world settings—particularly in the nonprofit sector—introduces additional design constraints related to accessibility, trust, and inclusion. 

Within the nonprofit and social-impact space, there is growing interest in deploying chatbots for information provision, advocacy and service navigation. Følstad et al.\ discuss “chatbots for social good’’ as a research and design agenda, identifying use cases in health, education, civic participation and humanitarian response, and highlighting issues of inclusion, trust, and governance \cite{folstad2018socialgood}. They note that many early systems are narrowly scoped, relying on scripted or intent-based flows.

Empirical work on health-related chatbots in low and middle income countries (LMICs) provides concrete examples. Munir et al.\ report on the feasibility of an “artificially intelligent vaccines chatbot’’ (Bablibot) in Pakistan, deployed via text messaging to provide immunization information to caregivers in low-income communities \cite{munir2021bablibot}. Their mixed-methods evaluation finds high user satisfaction and suggests that chatbots can address access barriers, especially for women. Building on this, Siddiqi et al.\ present a detailed mixed-methods study of Bablibot, describing its design using natural language processing (NLP), machine learning and a human-in-the-loop component, and show that local-language chatbots are feasible and acceptable in low-resource, low-literacy settings \cite{siddiqi2024bablibot}.

Beyond health, chatbots are increasingly used for fundraising and advocacy by nonprofit organizations, including for storytelling and campaign communication on platforms such as Facebook Messenger \cite{halperin2023storytelling}. However, many of these systems focus on a single function (e.g., advocacy narratives, fundraising journeys) and are not deeply integrated with internal programme databases or multi-department service structures.

\section{Pakistan Context and Local-Language Chatbots}
These challenges are further amplified in the Pakistani context, where linguistic diversity and resource constraints shape the design and deployment of conversational systems. Pakistan is an important context for AI-mediated services because it is a lower middle income, linguistically diverse country with high mobile and messaging usage.  Rich morphology, complex script, and cross-script usage are some of the unique NLP challenges faced by Urdu and Roman Urdu.  In their review of problems and difficulties in Urdu Natural Language Processing, Khan et al. point out the lack of annotated tools and resources and identify tasks like tokenisation, morphological analysis, POS tagging, parsing, and named-entity recognition as ongoing research issues \cite{khan2020urdunlp}.  The development of Urdu NLP resources and applications is also described in a more comprehensive survey of Urdu language processing \cite{anwar2016survey}.

Several recent projects have focused on Urdu and Roman Urdu in conversational AI. For instance, Kumar et al. developed SAATHI, a voice-enabled Urdu virtual assistant designed to help senior citizens in Pakistan who are “ageing in place” by providing news, reminders, and simple interactions in Urdu \cite{kumar2022saathi}. In another project, Mohiuddin et al. created a multilingual Urdu chatbot using the Rasa framework, which tackles the challenge of converting Roman Urdu input into Urdu intents \cite{mohiuddin2023multilingual}. Similarly, Hakeem et al. describe an AI-powered chatbot for Roman Urdu, combining GPT-style models with a RAG framework to improve handling of spelling variations and the lack of standardization \cite{hakeem2025romanurdu}.

At the same time, evaluation of LLMs on Urdu remains an active area. Kazi et al.\ evaluate large language models on Urdu–English question answering, showing significant performance gaps between high-resource and low-resource settings and underlining the need for Urdu-specific adaptation and datasets \cite{kazi2025llmqa}. These works collectively demonstrate that:
\begin{itemize}\setlength{\itemsep}{0pt}
    \item Urdu and Roman Urdu chatbots are technically feasible;
    \item local-language assistants can be deployed in low-resource Pakistani communities;
    \item there is still limited work on large, multi-domain assistants tightly integrated with a single NGO’s internal systems.
\end{itemize}

\section{Similar Work in NGO Service and Donor Assistants}
Building on these contextual considerations, prior work on NGO-oriented assistants can be grouped into distinct categories based on domain scope and system capabilities. In the specific domain of nonprofit and NGO assistants, the most relevant academic work can be grouped into three strands.

First, chatbots for social good and humanitarian applications. The CHI extended abstract on chatbots for social good explicitly frames chatbots as tools for health, education and civic engagement, and calls for systems that can scale while remaining inclusive and trustworthy \cite{folstad2018socialgood}. This provides a conceptual foundation for NGO-oriented assistants such as the Alkhidmat portal.

Second, health and immunization chatbots in Pakistan. Munir et al.\ and Siddiqi et al.\ both focus on Bablibot, a local-language chatbot for immunization information in Karachi \cite{munir2021bablibot,siddiqi2024bablibot}. Bablibot is tightly connected to immunization registries and demonstrates integration with digital health records, but it is single-domain (vaccines) and not a general NGO service portal.

Third, Urdu / Roman Urdu service chatbots and administrative bots. Mohiuddin et al.\ target multi-lingual service access via transliteration-based Urdu chatbots using Rasa \cite{mohiuddin2023multilingual}; Kumar et al.\ design SAATHI for elderly care information \cite{kumar2022saathi}; Hakeem et al.\ build a Roman Urdu chatbot using RAG for more open-ended conversational understanding \cite{hakeem2025romanurdu}. Maqbool et al.\ propose a sentence-level encoder–decoder architecture for an intelligent administrative chatbot for Roman Urdu, designed to support university administrative tasks \cite{maqbool2025adminbot}. These systems are structurally closer to an organizational assistant, but they are focused on administrative or general chat tasks rather than multi-department NGO services.

There is little evidence in the academic literature of a \emph{single} Pakistani NGO deploying a RAG-based assistant that simultaneously covers donations, healthcare navigation, blood services, welfare programmes, and volunteer opportunities across web and WhatsApp.

\section{Multimodal Inputs and Accessibility} 
Modern conversational assistants increasingly handle not just text but multiple input modalities, such as images and voice, significantly enhancing accessibility and interaction.

Recent work on Multimodal RAG (Retrieval-Augmented Generation) shows how chatbots can analyze and reason over images, tables, and text together to enrich responses. For instance, systems can perform OCR or caption images and retrieve relevant text or table data to answer user queries, reflecting an alignment with research on "multimodal" chatbots that combine vision and language for richer dialogue \cite{gulab2024multimodalrag}.

State-of-the-art models, such as GPT-4V ("GPT-4 with Vision"), now support direct image input, processing raw image bytes alongside text \cite{gulab2024multimodalrag}. OpenAI has added similar image and voice capabilities to ChatGPT, allowing users to "show ChatGPT what you’re talking about" and converse by voice \cite{openai2023multimodal}. An AI chatbot might use convolutional neural networks (CNNs) to parse objects and context from an uploaded photo \cite{ijaresm2024multimodal}.

In practical terms, this means our chatbot can accept photos of payment receipts, prescriptions, or donation screenshots and extract contextual information from them \cite{gulab2024multimodalrag}.

Voice interfaces are also becoming common for accessibility:
\begin{itemize}\setlength{\itemsep}{0pt}
    \item A speech-to-text (STT) or automatic speech recognition (ASR) component can transcribe user voice notes (with basic Urdu support) into the dialogue, allowing users who prefer speaking to interact with the bot  \cite{gulab2024multimodalrag}, \cite{ijaresm2024multimodal}.
    \item In a Pakistani NGO context, this means a donor could send a voice note in Urdu that is transcribed and answered via ASR, bridging literacy and connectivity gaps \cite{ijaresm2024multimodal}, \cite{openai2023multimodal}.
\end{itemize}

Integrating image and speech inputs enables the system to serve non-tech-savvy, low-literacy users, or those with disabilities, extending the chatbot beyond "plain text" to richer multimodal engagement \cite{gulab2024multimodalrag}. Such features make digital assistants more inclusive (e.g., by letting low-literacy users show receipts or speak questions) \cite{openai2023multimodal}.

These multimodal capabilities are rapidly dropping in cost and complexity as open-source vision models (e.g., CLIP, BLIP-2) and multimodal LLM APIs (e.g., GPT-4V, Claude Vision) become more accessible \cite{gulab2024multimodalrag}.

\section{Human-in-the-Loop and Escalation Management}
Industry best practices emphasize that chatbots should gracefully escalate complex or ambiguous queries to human agents to maintain user trust and ensure effectiveness \cite{freshworks2024analytics}. In hybrid AI support systems, escalation is typically triggered when the model’s confidence is low or when the user’s request falls outside the chatbot’s defined domain, such as medical emergencies or detailed calculations. These systems rely on monitoring intent confidence and transferring any ambiguous or sensitive query to staff to prevent misinformation and ensure reliability \cite{freshworks2024analytics}. Platforms such as WATI, a WhatsApp Business API provider, operationalize this principle by implementing hybrid workflows in which the AI bot “automatically answers common queries” but “routes more complex queries to your support teams/agents based on customer intent” \cite{wati2024aisupport}. In addition, industry vendors stress that human agents must retain complete control to manage AI conversations manually when required, ensuring operational oversight, accountability, and continuity in user experience.

\subsection{Uncertainty Estimation and Confidence Scoring in RAG Systems}
A key challenge in implementing human-in-the-loop systems is determining \emph{when} escalation should occur. This requires explicit estimation of model confidence and uncertainty, which has become an active area of research in LLM-based systems. Reliable deployment of RAG-based assistants requires explicit modeling of uncertainty to decide when to answer, abstain, or escalate. Recent work categorizes confidence estimation methods into four classes: token-level, self-verbalized, semantic similarity-based, and mechanistic approaches \cite{acm2024uncertainty}.

\textbf{Token-level methods} leverage the probabilistic nature of LLMs. Since models generate tokens sequentially with associated probabilities, confidence can be derived from log-probabilities. Common techniques include:
\begin{itemize}\setlength{\itemsep}{0pt}
    \item \textit{Average log-probability} and \textit{maximum token uncertainty} for sequence-level confidence;
    \item \textit{Perplexity}, where higher values indicate greater uncertainty due to a flatter probability distribution;
    \item \textit{Entropy-based scoring}, where low entropy reflects confident (peaked) token distributions and high entropy indicates uncertainty \cite{acm2024uncertainty}.
\end{itemize}

Recent applied approaches propose weighted confidence formulations that emphasize high-certainty tokens. For example, combining the joint probability of top-ranked tokens with overall sequence probability (e.g., 70/30 weighting) has been shown to stabilize confidence estimates in generation tasks \cite{thomas2024confidence}.

\textbf{Sampling-based methods}, such as SelfCheckGPT, estimate confidence by generating multiple responses and measuring consistency across them. Higher agreement across sampled outputs implies higher confidence, making this approach suitable for black-box models where logits are unavailable, albeit at higher computational cost \cite{capgemini2024uncertainty}.

\textbf{Self-verbalized confidence} asks the model to report its own certainty. While simple, this approach is prone to overconfidence and is typically used in combination with other signals \cite{capgemini2024uncertainty}.

\textbf{Retrieval-aware confidence} is also critical in RAG systems. A common proxy is the semantic similarity between the query and retrieved documents (e.g., cosine similarity), aggregated across top-$k$ chunks to estimate retrieval reliability \cite{geeksforgeeks2024ragmetrics}. This complements generation-side confidence.

\textbf{Threshold selection} plays a central role in operationalizing confidence. Industry and research suggest tuning thresholds empirically on validation data to optimize correctness, often by sweeping values between 0 and 1 and selecting those that maximize accuracy for high-confidence responses \cite{mit2022overconfidence,cyara2023threshold}. These thresholds directly inform fallback strategies such as abstention or human escalation.

\textbf{Implication for system design:} In the proposed system, confidence is modeled as a hybrid signal combining (i) token-level generation confidence, (ii) retrieval similarity scores, and (iii) optional sampling-based validation for critical queries. This composite score governs response acceptance, rejection (“I don’t have enough information”), and routing to human agents, aligning with best practices for trustworthy AI deployment.


\section{Staff Coordination and Agent Dashboard}
The operationalization of human-in-the-loop handoff is facilitated through shared agent dashboards, often referred to as “team inboxes,” as exemplified by platforms such as WATI \cite{wati2024aisupport}. These dashboards are designed to enable seamless coordination among staff when managing escalated queries. Each escalated chat is presented as a ticket, allowing agents to access the complete conversation history and respond directly through the same channel, including bilingual responses where required \cite{wati2024aisupport}. Domain-based routing capabilities are commonly integrated, ensuring that specific queries, such as those related to donations, are directed to the appropriate team (e.g., fundraising) and can be dynamically reassigned among specialists. The inclusion of pre-approved templates supports standardized and high-quality responses, particularly within platform-specific constraints such as WhatsApp’s 24-hour reply window, while the logging of agent actions provides traceability and auditability. Furthermore, these systems are designed to maintain continuous availability and multilingual support, enabling consistent handoff and follow-up interactions irrespective of whether the user communicates in English, Urdu, or Roman Urdu. Vendors highlight that such implementations provide “24/7 automated support that responds naturally and accurately in your customers’ preferred language” \cite{wati2024aisupport}, thereby ensuring both operational efficiency and user satisfaction.

\section{Analytics and Admin Tools}
Effective deployments of organizational chatbots require comprehensive analytics and content management to guide continuous improvement, maintain accuracy, and ensure transparency over time. Performance monitoring is typically conducted through analytics dashboards that report Key Performance Indicators (KPIs) such as query volume, total queries handled, and deflection rate, indicating the percentage of interactions resolved by the AI versus those escalated to human agents. Industry guides emphasize escalation or human-takeover rate as a vital KPI, noting that low takeover rates suggest strong autonomous handling, whereas spikes may signal knowledge gaps or safety concerns \cite{freshworks2024analytics}. Additional time-based metrics commonly tracked include average response time and session duration, and systems routinely collect user satisfaction data, including CSAT (thumbs-up/down), to assess perceived quality \cite{freshworks2024analytics}. Real-time dashboards also chart usage by domain (e.g., donations or health queries) and peak hours, enabling administrators to interpret logs, review conversation transcripts, and incorporate user feedback to iteratively improve the bot, refine intents or flows, and update the RAG knowledge base \cite{freshworks2024analytics}. The literature consistently underscores that chatbots must function as “data-driven” tools \cite{freshworks2024analytics}.

A human-facing administrative interface complements analytics by enabling operational control and content management. Staff require the ability to manually update program information, such as donation campaigns and hospital details, and maintain the knowledge repository. In the context of a Pakistani NGO, such an interface would also allow editing of Urdu translations, defining personas, and monitoring escalation tickets in real time. Administrative configuration capabilities typically include adjusting fallback messages, modifying quick-reply templates, and setting content moderation rules. Systems comparable to enterprise chat platforms incorporate “team inbox” operational views alongside update forms for knowledge base maintenance \cite{wati2024aisupport}. Together, automated analytics and manual administrative controls align with recognized best practices in RAG-based assistant deployment, ensuring sustained accuracy, transparency, and ongoing system improvement.

\section{Cloud Deployment and Scalability}
To accommodate large user volumes, ensure high availability, and facilitate rapid growth, conversational platforms are typically deployed on cloud infrastructure with containerized architectures. Many organizations, including non-governmental organizations (NGOs), leverage Platform-as-a-Service (PaaS) solutions such as Heroku, AWS, and Google Cloud, which provide autoscaling capabilities, managed databases, and high availability \cite{heroku2024chatbots}, \cite{aws2024marketplace}. Salesforce’s Live Agent chatbot, for instance, is deployed on Heroku, utilizing managed Postgres databases and benefiting from dedicated operational staff responsible for backups, updates, and outage management, thereby reducing the operational burden on developers \cite{heroku2024chatbots}. Heroku advertises its PaaS offering as providing “fully managed infrastructure... with built-in autoscaling and detailed performance metrics” in a “compliant environment.”  

The use of PaaS or containerized cloud clusters ensures both high availability and seamless horizontal scalability, enabling the system to accommodate hundreds of simultaneous users and remain resilient during peak usage periods \cite{heroku2024chatbots}. Managed cloud services further enhance operational efficiency by providing automated backups, secure key management, and CI/CD pipelines for version control, along with real-time performance monitoring \cite{heroku2024chatbots}, \cite{aws2024marketplace}. Technical components of such deployments often include Docker containers, and services like Elasticsearch or pgvector for Retrieval-Augmented Generation (RAG) indexing, which support maintainable and responsive chatbot systems. Leveraging these cloud-based practices enables platforms to reliably support large-scale user interactions while maintaining uptime exceeding 99\%.

\section{Security, Privacy, and Compliance}
For NGOs deploying conversational agents, safeguarding user data is of paramount importance and must align with both industry standards and relevant legal frameworks. The architecture of such systems should integrate robust encryption mechanisms, ensuring that all communications, particularly those containing sensitive personal or medical information, are secured in transit and at rest. For web interactions, HTTPS/TLS protocols should be applied, while WhatsApp-based interactions benefit from end-to-end Signal encryption inherent to the platform \cite{whatsapp2024trust}.  

Beyond technical security, adherence to global compliance standards is critical. The WhatsApp Business Platform holds certifications including SOC 2 and explicitly supports GDPR compliance, providing a framework for lawful data processing \cite{whatsapp2024trust}. In addition, organizations operating in Pakistan should comply with emerging national regulations, such as the draft Personal Data Protection Bill (2023), which mirrors GDPR by mandating user consent, data minimization, and breach accountability \cite{chambers2025pakistan}.  

Data minimization policies are central to ethical and legal compliance. Best practices include obtaining explicit user consent \cite{whatsapp2024trust}, \cite{chambers2025pakistan}, retaining only transient session logs, and deleting chat histories after completion of interactions \cite{whatsapp2024trust}. Personal identifiers should be anonymized, and sensitive personally identifiable information (PII) should not be stored beyond operational necessity \cite{chambers2025pakistan}.  

In conclusion, a secure and compliant chatbot deployment integrates the built-in end-to-end encryption and certifications of the WhatsApp Business API while implementing internal policies for data minimization and local regulatory compliance. This combination ensures both technical security and adherence to legal and ethical standards, providing a reliable framework for NGO services \cite{whatsapp2024trust}, \cite{chambers2025pakistan}.

\section{Differences Between Existing Work and the Proposed System}
Synthesizing the above, the proposed Public Chat Portal for Alkhidmat differs from existing work in several ways.

\subsection{Multi-Domain NGO Scope Versus Single-Domain Chatbots}
Existing Pakistani chatbots in the literature typically focus on one domain: immunization (Bablibot) \cite{munir2021bablibot,siddiqi2024bablibot}, elderly care information (SAATHI) \cite{kumar2022saathi}, administrative tasks in education \cite{maqbool2025adminbot}, or general Roman Urdu conversational understanding \cite{hakeem2025romanurdu}. In contrast, the Alkhidmat assistant is explicitly designed as a multi-department organizational portal, covering donations (e.g., Zakat, Sadaqah), healthcare facilities, blood services (donation and requests), other general facilities.

\subsection{Deep RAG Over NGO-Specific Knowledge}
Most NGO-relevant chatbots described in the literature are intent or flow based, with relatively shallow use of organizational documents or databases \cite{folstad2018socialgood,mohiuddin2023multilingual}. By contrast, the proposed system follows the RAG paradigm discussed by Lewis et al.\ and later surveys \cite{lewis2020rag,gao2023ragsurvey,fan2024ragllm}, building a structured knowledge base from Alkhidmat’s internal SOPs, programme descriptions, facility lists, and operational data, and using retrieval to ground responses. This is closer to enterprise RAG assistants but applied in a nonprofit setting.

\subsection{Localization to Urdu, Roman Urdu, and Pakistani NGO Practices}
While prior work has shown that Urdu and Roman Urdu chatbots are feasible \cite{khan2020urdunlp,kumar2022saathi,mohiuddin2023multilingual,hakeem2025romanurdu}, these systems are typically either domain-specific (health, elderly care) or experimental. The Alkhidmat assistant aims to:
\begin{itemize}\setlength{\itemsep}{0pt}
    \item support English, Urdu, and basic Roman Urdu input;
    \item incorporate programme names and concepts specific to a large Pakistani NGO (e.g., Mawakhat, Orphan Care, Bano Qabil);
    \item operate on channels (web and WhatsApp Business) that match actual user behaviour in Pakistan.
\end{itemize}
This combination of linguistic and institutional localization is not extensively covered in existing academic work.

\subsection{Integration With Operational Systems and Human Escalation}
The Bablibot work demonstrates integration with immunization registries and a human-in-the-loop support model \cite{siddiqi2024bablibot,munir2021bablibot}. The proposed Alkhidmat portal extends this logic beyond health to other NGO operations, for example:
\begin{itemize}\setlength{\itemsep}{0pt}
    \item a blood-donation microservice that can check donor eligibility criteria and facility locations;
    \item dashboards where staff can view escalated conversations with context and respond directly;
    \item analytics on domain-wise query volumes, response quality, and escalation rates.
\end{itemize}
This reflects best practices in RAG-based enterprise assistants (monitoring, feedback loops) but adapted for NGO governance.
Escalation decisions are typically driven by confidence thresholds derived from uncertainty estimation. Queries with low generation confidence or weak retrieval grounding are automatically routed to human agents, ensuring reliability and preventing hallucinated responses \cite{mit2022overconfidence,cyara2023threshold}.

\section{Conclusion}
The literature establishes Retrieval-Augmented Generation as a robust framework for building grounded and reliable conversational systems. At the same time, research in conversational AI emphasizes the importance of multi-turn interaction, persona adaptation, and user trust. In nonprofit contexts, chatbot deployments demonstrate feasibility and impact, particularly in LMIC settings, but remain largely narrow in scope and function.